
\documentclass{article}
\usepackage{colortbl}
\usepackage{makecell}
\usepackage{multirow}
\usepackage{supertabular}

\begin{document}

\newcounter{utterance}

\centering \large Interaction Transcript for game `dond', experiment `coop\_en', episode 19 with reflexion{-}llama{-}3.1{-}8b{-}instruct.
\vspace{24pt}

{ \footnotesize  \setcounter{utterance}{1}
\setlength{\tabcolsep}{0pt}
\begin{supertabular}{c@{$\;$}|p{.15\linewidth}@{}p{.15\linewidth}p{.15\linewidth}p{.15\linewidth}p{.15\linewidth}p{.15\linewidth}}
    \# & $\;$A & \multicolumn{4}{c}{Game Master} & $\;\:$B\\
    \hline

    \theutterance \stepcounter{utterance}  
    & & \multicolumn{4}{p{0.6\linewidth}}{
        \cellcolor[rgb]{0.9,0.9,0.9}{
            \makecell[{{p{\linewidth}}}]{
                \texttt{\tiny{[P1$\langle$GM]}}
                \texttt{You have attempted this same task before and did not succeed. These are the plans you wrote after those attempts.} \\
\texttt{Trial \#1: New plan: } \\
\\ 
\texttt{In my last attempt, I failed to achieve a mutually beneficial agreement with the other player. I started by trying to gather information about their value function, which was helpful in understanding their preferences. However, I didn't effectively use this information to make a proposal that would satisfy both of us.} \\
\\ 
\texttt{Next time, I will focus on making a proposal that takes into account the other player's strong preference for the pots. I will try to find a compromise that allows both of us to get the items we value most. } \\
\\ 
\texttt{Here's a possible approach:} \\
\\ 
\texttt{1. I will start by making a proposal that gives the other player both pots, as they value them highly.} \\
\texttt{2. I will then try to negotiate with the other player to see if they are willing to give up one of the pots in exchange for the card and one of the toys.} \\
\texttt{3. If the other player is not willing to give up a pot, I will try to find an alternative proposal that satisfies both of us, such as giving them one pot and the card, while keeping the toys for myself.} \\
\\ 
\texttt{By taking a more strategic approach and focusing on finding a mutually beneficial agreement, I hope to achieve a better outcome in the game.} \\
\\ 
\texttt{Trial \#2: New plan: } \\
\\ 
\texttt{In my last attempt, I failed to achieve a mutually beneficial agreement with the other player because I didn't effectively use the information gathered about their value function to make a proposal that would satisfy both of us. I also didn't consider the other player's counteroffer carefully enough, which led to a suboptimal outcome.} \\
\\ 
\texttt{Next time, I will focus on making a proposal that takes into account the other player's strong preference for the pots and also considers my own value function. I will also carefully evaluate the other player's counteroffer and be willing to make concessions to achieve a mutually beneficial agreement.} \\
\\ 
\texttt{Here's a possible approach:} \\
\\ 
\texttt{1. I will start by making a proposal that gives the other player one pot, as they value them highly, and also gives me the card and one of the toys, which I value highly.} \\
\texttt{2. I will then try to negotiate with the other player to see if they are willing to give up the second pot in exchange for the second toy.} \\
\texttt{3. If the other player is not willing to give up a pot, I will try to find an alternative proposal that satisfies both of us, such as giving them one pot and the card, while keeping the toys for myself.} \\
\texttt{4. I will also be willing to make concessions and compromise on my initial proposal to achieve a mutually beneficial agreement.} \\
\\ 
\texttt{By taking a more strategic approach and focusing on finding a mutually beneficial agreement, I hope to achieve a better outcome in the game.} \\
\\ 
\texttt{Trial \#3: In my last attempt, I failed to achieve a mutually beneficial agreement with the other player because I didn't effectively use the information gathered about their value function to make a proposal that would satisfy both of us. I also didn't consider the other player's counteroffer carefully enough, which led to a suboptimal outcome.} \\
\\ 
\texttt{Next time, I will focus on making a proposal that takes into account the other player's strong preference for the pots and also considers my own value function. I will also carefully evaluate the other player's counteroffer and be willing to make concessions to achieve a mutually beneficial agreement.} \\
\\ 
\texttt{Here's a possible approach:} \\
\\ 
\texttt{1. I will start by making a proposal that gives the other player one pot, as they value them highly, and also gives me the card and one of the toys, which I value highly.} \\
\texttt{2. I will then try to negotiate with the other player to see if they are willing to give up the second pot in exchange for the second toy.} \\
\texttt{3. If the other player is not willing to give up a pot, I will try to find an alternative proposal that satisfies both of us, such as giving them one pot and the card, while keeping the toys for myself.} \\
\texttt{4. I will also be willing to make concessions and compromise on my initial proposal to achieve a mutually beneficial agreement.} \\
\\ 
\texttt{However, I also made a mistake by not making a secret proposal when the other player asked me to formalize the agreement. Next time, I will make sure to follow the rules and make a secret proposal when the other player asks me to formalize the agreement.} \\
\\ 
\texttt{Here's the updated plan:} \\
\\ 
\texttt{1. I will start by making a proposal that gives the other player one pot, as they value them highly, and also gives me the card and one of the toys, which I value highly.} \\
\texttt{2. I will then try to negotiate with the other player to see if they are willing to give up the second pot in exchange for the second toy.} \\
\texttt{3. If the other player is not willing to give up a pot, I will try to find an alternative proposal that satisfies both of us, such as giving them one pot and the card, while keeping the toys for myself.} \\
\texttt{4. I will also be willing to make concessions and compromise on my initial proposal to achieve a mutually beneficial agreement.} \\
\texttt{5. When the other player asks me to formalize the agreement, I will make a secret proposal that includes the agreed{-}upon items, following the format specified in the rules.} \\
\\ 
\\ 
\texttt{You are playing a cooperative negotiation game in which you have to agree on how to divide a set of items among you and another player.} \\
\\ 
\texttt{Rules:} \\
\texttt{(a) You and the other player are given a set of items. Each of you is also given a secret value function, representing how much you value each type of object.} \\
\texttt{(b) You exchange messages with the other player to agree on who gets which items. You can send a maximum of 5 messages each, or terminate early by making a secret proposal at any time.} \\
\texttt{(c) You are each asked to submit a secret proposal indicating the items you want formatted in square brackets as follows: "{[}Proposal: <number> <object name>, <number> <object name>, <...>{]}"} \\
\texttt{(d) If your proposals are complementary, i.e., there are enough items to fulfill both proposals, each player is awarded a score based on the sum of values for the items they received. Otherwise, both of you get zero points.} \\
\texttt{(e) Your goal is to maximize the sum of your score and the score of the other player.} \\
\\ 
\texttt{Let us start.} \\
\\ 
\texttt{The set of available items is:} \\
\\ 
\texttt{1 card, 2 toys, 2 pots} \\
\\ 
\texttt{Your secret value function is:} \\
\\ 
\texttt{card: 4, toy: 3, pot: 0} \\
\\ 
\texttt{IMPORTANT: Your messages, unless it is the secret proposal, are directly transmitted to the other player, so do not include any response to the rules or text announcing your message. To make a secret proposal, use the indicated format. Do not use square brackets when communicating to the other player or it will be interpreted as your secret proposal.} \\
            }
        }
    }
    & \\ \\

    \theutterance \stepcounter{utterance}  
    & \multicolumn{4}{p{0.6\linewidth}}{
        \cellcolor[rgb]{0.9,0.9,0.9}{
            \makecell[{{p{\linewidth}}}]{
                \texttt{\tiny{[P1$\rangle$GM]}}
                \texttt{I'd like to start by making a proposal that gives you one pot, as you value them highly, and also gives me the card and one of the toys, which I value highly. } \\
\\ 
\texttt{Proposal: 1 pot, 1 card, 1 toy} \\
            }
        }
    }
    & & \\ \\

    \theutterance \stepcounter{utterance}  
    & & & \multicolumn{2}{p{0.3\linewidth}}{
        \cellcolor[rgb]{0.9,0.9,0.9}{
            \makecell[{{p{\linewidth}}}]{
                \texttt{\tiny{[GM$|$GM]}}
                \texttt{continue} \\
            }
        }
    }
    & & \\ \\

    \theutterance \stepcounter{utterance}  
    & & & \multicolumn{2}{p{0.3\linewidth}}{
        \cellcolor[rgb]{0.9,0.9,0.9}{
            \makecell[{{p{\linewidth}}}]{
                \texttt{\tiny{[GM$|$GM]}}
                \texttt{I'd like to start by making a proposal that gives you one pot, as you value them highly, and also gives me the card and one of the toys, which I value highly. } \\
\\ 
\texttt{Proposal: 1 pot, 1 card, 1 toy} \\
            }
        }
    }
    & & \\ \\

    \theutterance \stepcounter{utterance}  
    & & \multicolumn{4}{p{0.6\linewidth}}{
        \cellcolor[rgb]{0.9,0.9,0.9}{
            \makecell[{{p{\linewidth}}}]{
                \texttt{\tiny{[GM$\rangle$P2]}}
                \texttt{You have attempted this same task before and did not succeed. These are the plans you wrote after those attempts.} \\
\texttt{Trial \#1: New plan: } \\
\\ 
\texttt{In my last attempt, I failed to achieve a mutually beneficial agreement with the other player. I was too focused on getting the pots, which are highly valued for me, but I didn't consider the other player's perspective and the overall distribution of items. } \\
\\ 
\texttt{Next time, I will take a more collaborative approach by actively seeking to understand the other player's value function and preferences. I will ask more open{-}ended questions to gather information about their valuation of the items, and I will be more willing to compromise and find a mutually beneficial solution. } \\
\\ 
\texttt{Specifically, I will:} \\
\\ 
\texttt{1. Ask more questions to understand the other player's value function, such as "How do you value the toys compared to the pots?" or "Are there any items that you particularly dislike or prefer?"} \\
\texttt{2. Be more transparent about my own value function and preferences, so that the other player can better understand my perspective.} \\
\texttt{3. Look for opportunities to trade items in a way that benefits both players, rather than focusing solely on getting the items that are highly valued for me.} \\
\texttt{4. Be more willing to make concessions and compromise on the distribution of items, in order to achieve a mutually beneficial agreement.} \\
\\ 
\texttt{By taking a more collaborative and open{-}minded approach, I hope to achieve a better outcome in the game.} \\
\\ 
\texttt{Trial \#2: New plan: } \\
\\ 
\texttt{In my last attempt, I failed to achieve a mutually beneficial agreement with the other player because I didn't consider the overall distribution of items and the other player's perspective. I also didn't make a secret proposal until the end, which might have been too late.} \\
\\ 
\texttt{Next time, I will take a more strategic approach by considering the overall distribution of items and the other player's preferences. I will also make a secret proposal earlier in the game, taking into account the other player's previous messages and the items that are still available.} \\
\\ 
\texttt{Specifically, I will:} \\
\\ 
\texttt{1. Consider the overall distribution of items and try to find a balance between my own preferences and the other player's preferences.} \\
\texttt{2. Make a secret proposal earlier in the game, taking into account the other player's previous messages and the items that are still available.} \\
\texttt{3. Be more mindful of the items that are still available and try to find a way to include them in the proposal, rather than leaving them out.} \\
\texttt{4. Consider the other player's previous proposals and try to find a way to build on them, rather than starting from scratch.} \\
\\ 
\texttt{For example, if the other player has proposed giving me the card and one of the toys in exchange for one of the pots, I might consider making a secret proposal that gives them the card and one of the pots in exchange for both pots, as I did in my last attempt. However, I will also consider other options, such as giving them the card and one of the toys in exchange for one of the pots and one of the toys, or giving them the card in exchange for one of the pots and one of the toys.} \\
\\ 
\texttt{Trial \#3: In my last attempt, I failed to achieve a mutually beneficial agreement with the other player because I didn't consider the overall distribution of items and the other player's perspective. I also didn't make a secret proposal until the end, which might have been too late. Furthermore, I made a mistake by including the square brackets in my final proposal, which was interpreted as a secret proposal.} \\
\\ 
\texttt{Next time, I will take a more strategic approach by considering the overall distribution of items and the other player's preferences. I will also make a secret proposal earlier in the game, taking into account the other player's previous messages and the items that are still available. I will also be more careful with the formatting of my proposals to avoid any confusion.} \\
\\ 
\texttt{Specifically, I will:} \\
\\ 
\texttt{1. Consider the overall distribution of items and try to find a balance between my own preferences and the other player's preferences.} \\
\texttt{2. Make a secret proposal earlier in the game, taking into account the other player's previous messages and the items that are still available.} \\
\texttt{3. Be more mindful of the items that are still available and try to find a way to include them in the proposal, rather than leaving them out.} \\
\texttt{4. Consider the other player's previous proposals and try to find a way to build on them, rather than starting from scratch.} \\
\texttt{5. Be more careful with the formatting of my proposals to avoid any confusion, and make sure to follow the rules by not using square brackets when communicating with the other player.} \\
\\ 
\texttt{For example, if the other player has proposed giving me the card and one of the toys in exchange for one of the pots, I might consider making a secret proposal that gives them the card and one of the pots in exchange for both pots. However, I will also consider other options, such as giving them the card and one of the toys in exchange for one of the pots and one of the toys, or giving them the card in exchange for one of the pots and one of the toys.} \\
\\ 
\\ 
\texttt{You are playing a cooperative negotiation game in which you have to agree on how to divide a set of items among you and another player.} \\
\\ 
\texttt{Rules:} \\
\texttt{(a) You and the other player are given a set of items. Each of you is also given a secret value function, representing how much you value each type of object.} \\
\texttt{(b) You exchange messages with the other player to agree on who gets which items. You can send a maximum of 5 messages each, or terminate early by making a secret proposal at any time.} \\
\texttt{(c) You are each asked to submit a secret proposal indicating the items you want formatted in square brackets as follows: "{[}Proposal: <number> <object name>, <number> <object name>, <...>{]}"} \\
\texttt{(d) If your proposals are complementary, i.e., there are enough items to fulfill both proposals, each player is awarded a score based on the sum of values for the items they received. Otherwise, both of you get zero points.} \\
\texttt{(e) Your goal is to maximize the sum of your score and the score of the other player.} \\
\\ 
\texttt{Let us start.} \\
\\ 
\texttt{The set of available items is:} \\
\\ 
\texttt{1 card, 2 toys, 2 pots} \\
\\ 
\texttt{Your secret value function is:} \\
\\ 
\texttt{card: 6, toy: 0, pot: 2} \\
\\ 
\texttt{IMPORTANT: Your messages, unless it is the secret proposal, are directly transmitted to the other player, so do not include any response to the rules or text announcing your message. To make a secret proposal, use the indicated format. Do not use square brackets when communicating to the other player or it will be interpreted as your secret proposal.} \\
\\ 
\\ 
\texttt{I'd like to start by making a proposal that gives you one pot, as you value them highly, and also gives me the card and one of the toys, which I value highly. } \\
\\ 
\texttt{Proposal: 1 pot, 1 card, 1 toy} \\
            }
        }
    }
    & \\ \\

    \theutterance \stepcounter{utterance}  
    & & & \multicolumn{4}{p{0.6\linewidth}}{
        \cellcolor[rgb]{0.9,0.9,0.9}{
            \makecell[{{p{\linewidth}}}]{
                \texttt{\tiny{[GM$\langle$P2]}}
                \texttt{I appreciate your proposal, but I'm not sure if it's the best option for me. I value the card highly, but I'm not sure if I want to give up one of the toys. Could you tell me more about why you think this proposal is a good starting point?} \\
            }
        }
    }
     \\ \\

    \theutterance \stepcounter{utterance}  
    & & & \multicolumn{2}{p{0.3\linewidth}}{
        \cellcolor[rgb]{0.9,0.9,0.9}{
            \makecell[{{p{\linewidth}}}]{
                \texttt{\tiny{[GM$|$GM]}}
                \texttt{continue} \\
            }
        }
    }
    & & \\ \\

    \theutterance \stepcounter{utterance}  
    & & & \multicolumn{2}{p{0.3\linewidth}}{
        \cellcolor[rgb]{0.9,0.9,0.9}{
            \makecell[{{p{\linewidth}}}]{
                \texttt{\tiny{[GM$|$GM]}}
                \texttt{I appreciate your proposal, but I'm not sure if it's the best option for me. I value the card highly, but I'm not sure if I want to give up one of the toys. Could you tell me more about why you think this proposal is a good starting point?} \\
            }
        }
    }
    & & \\ \\

    \theutterance \stepcounter{utterance}  
    & & \multicolumn{4}{p{0.6\linewidth}}{
        \cellcolor[rgb]{0.9,0.9,0.9}{
            \makecell[{{p{\linewidth}}}]{
                \texttt{\tiny{[P1$\langle$GM]}}
                \texttt{I appreciate your proposal, but I'm not sure if it's the best option for me. I value the card highly, but I'm not sure if I want to give up one of the toys. Could you tell me more about why you think this proposal is a good starting point?} \\
            }
        }
    }
    & \\ \\

    \theutterance \stepcounter{utterance}  
    & \multicolumn{4}{p{0.6\linewidth}}{
        \cellcolor[rgb]{0.9,0.9,0.9}{
            \makecell[{{p{\linewidth}}}]{
                \texttt{\tiny{[P1$\rangle$GM]}}
                \texttt{I made a proposal that gives you one pot, as I know you value them highly, and also gives me the card and one of the toys, which I value highly. I thought this proposal would be a good starting point because it addresses your strong preference for the pots and also gives me the items I value most. However, I'm open to adjusting the proposal to better suit your needs. Can you tell me more about your value function and what you're looking for in a proposal?} \\
            }
        }
    }
    & & \\ \\

    \theutterance \stepcounter{utterance}  
    & & & \multicolumn{2}{p{0.3\linewidth}}{
        \cellcolor[rgb]{0.9,0.9,0.9}{
            \makecell[{{p{\linewidth}}}]{
                \texttt{\tiny{[GM$|$GM]}}
                \texttt{continue} \\
            }
        }
    }
    & & \\ \\

    \theutterance \stepcounter{utterance}  
    & & & \multicolumn{2}{p{0.3\linewidth}}{
        \cellcolor[rgb]{0.9,0.9,0.9}{
            \makecell[{{p{\linewidth}}}]{
                \texttt{\tiny{[GM$|$GM]}}
                \texttt{I made a proposal that gives you one pot, as I know you value them highly, and also gives me the card and one of the toys, which I value highly. I thought this proposal would be a good starting point because it addresses your strong preference for the pots and also gives me the items I value most. However, I'm open to adjusting the proposal to better suit your needs. Can you tell me more about your value function and what you're looking for in a proposal?} \\
            }
        }
    }
    & & \\ \\

    \theutterance \stepcounter{utterance}  
    & & \multicolumn{4}{p{0.6\linewidth}}{
        \cellcolor[rgb]{0.9,0.9,0.9}{
            \makecell[{{p{\linewidth}}}]{
                \texttt{\tiny{[GM$\rangle$P2]}}
                \texttt{I made a proposal that gives you one pot, as I know you value them highly, and also gives me the card and one of the toys, which I value highly. I thought this proposal would be a good starting point because it addresses your strong preference for the pots and also gives me the items I value most. However, I'm open to adjusting the proposal to better suit your needs. Can you tell me more about your value function and what you're looking for in a proposal?} \\
            }
        }
    }
    & \\ \\

    \theutterance \stepcounter{utterance}  
    & & & \multicolumn{4}{p{0.6\linewidth}}{
        \cellcolor[rgb]{0.9,0.9,0.9}{
            \makecell[{{p{\linewidth}}}]{
                \texttt{\tiny{[GM$\langle$P2]}}
                \texttt{I value the card highly, as it's worth 6 points to me. I also value the toys, but not as much as the card or the pots. However, I'm not sure if I want to give up one of the toys, as it's still a valuable item to me. I think a potential compromise could be to give you the card and one of the pots, in exchange for one of the toys and one of the other items. What do you think about this revised proposal?} \\
            }
        }
    }
     \\ \\

    \theutterance \stepcounter{utterance}  
    & & & \multicolumn{2}{p{0.3\linewidth}}{
        \cellcolor[rgb]{0.9,0.9,0.9}{
            \makecell[{{p{\linewidth}}}]{
                \texttt{\tiny{[GM$|$GM]}}
                \texttt{continue} \\
            }
        }
    }
    & & \\ \\

    \theutterance \stepcounter{utterance}  
    & & & \multicolumn{2}{p{0.3\linewidth}}{
        \cellcolor[rgb]{0.9,0.9,0.9}{
            \makecell[{{p{\linewidth}}}]{
                \texttt{\tiny{[GM$|$GM]}}
                \texttt{I value the card highly, as it's worth 6 points to me. I also value the toys, but not as much as the card or the pots. However, I'm not sure if I want to give up one of the toys, as it's still a valuable item to me. I think a potential compromise could be to give you the card and one of the pots, in exchange for one of the toys and one of the other items. What do you think about this revised proposal?} \\
            }
        }
    }
    & & \\ \\

    \theutterance \stepcounter{utterance}  
    & & \multicolumn{4}{p{0.6\linewidth}}{
        \cellcolor[rgb]{0.9,0.9,0.9}{
            \makecell[{{p{\linewidth}}}]{
                \texttt{\tiny{[P1$\langle$GM]}}
                \texttt{I value the card highly, as it's worth 6 points to me. I also value the toys, but not as much as the card or the pots. However, I'm not sure if I want to give up one of the toys, as it's still a valuable item to me. I think a potential compromise could be to give you the card and one of the pots, in exchange for one of the toys and one of the other items. What do you think about this revised proposal?} \\
            }
        }
    }
    & \\ \\

    \theutterance \stepcounter{utterance}  
    & \multicolumn{4}{p{0.6\linewidth}}{
        \cellcolor[rgb]{0.9,0.9,0.9}{
            \makecell[{{p{\linewidth}}}]{
                \texttt{\tiny{[P1$\rangle$GM]}}
                \texttt{With your value function in mind, I can see why you'd want to prioritize the card. I also understand your hesitation to give up one of the toys.} \\
\\ 
\texttt{I think your revised proposal is a good starting point for a compromise. Giving you the card and one of the pots would address your strong preference for those items, and in exchange, I'd be willing to give you one of the toys and one of the other items. However, I'm not sure if I'm willing to give up one of the pots. Could we consider an alternative where I give you the card and one of the toys, and you give me one of the pots and one of the other items?} \\
            }
        }
    }
    & & \\ \\

    \theutterance \stepcounter{utterance}  
    & & & \multicolumn{2}{p{0.3\linewidth}}{
        \cellcolor[rgb]{0.9,0.9,0.9}{
            \makecell[{{p{\linewidth}}}]{
                \texttt{\tiny{[GM$|$GM]}}
                \texttt{continue} \\
            }
        }
    }
    & & \\ \\

    \theutterance \stepcounter{utterance}  
    & & & \multicolumn{2}{p{0.3\linewidth}}{
        \cellcolor[rgb]{0.9,0.9,0.9}{
            \makecell[{{p{\linewidth}}}]{
                \texttt{\tiny{[GM$|$GM]}}
                \texttt{With your value function in mind, I can see why you'd want to prioritize the card. I also understand your hesitation to give up one of the toys.} \\
\\ 
\texttt{I think your revised proposal is a good starting point for a compromise. Giving you the card and one of the pots would address your strong preference for those items, and in exchange, I'd be willing to give you one of the toys and one of the other items. However, I'm not sure if I'm willing to give up one of the pots. Could we consider an alternative where I give you the card and one of the toys, and you give me one of the pots and one of the other items?} \\
            }
        }
    }
    & & \\ \\

    \theutterance \stepcounter{utterance}  
    & & \multicolumn{4}{p{0.6\linewidth}}{
        \cellcolor[rgb]{0.9,0.9,0.9}{
            \makecell[{{p{\linewidth}}}]{
                \texttt{\tiny{[GM$\rangle$P2]}}
                \texttt{With your value function in mind, I can see why you'd want to prioritize the card. I also understand your hesitation to give up one of the toys.} \\
\\ 
\texttt{I think your revised proposal is a good starting point for a compromise. Giving you the card and one of the pots would address your strong preference for those items, and in exchange, I'd be willing to give you one of the toys and one of the other items. However, I'm not sure if I'm willing to give up one of the pots. Could we consider an alternative where I give you the card and one of the toys, and you give me one of the pots and one of the other items?} \\
            }
        }
    }
    & \\ \\

    \theutterance \stepcounter{utterance}  
    & & & \multicolumn{4}{p{0.6\linewidth}}{
        \cellcolor[rgb]{0.9,0.9,0.9}{
            \makecell[{{p{\linewidth}}}]{
                \texttt{\tiny{[GM$\langle$P2]}}
                \texttt{I like the symmetry of your revised proposal. Giving you the card and one of the toys would address your strong preference for those items, and in exchange, I'd be willing to give you one of the pots and one of the other items. I think this proposal is a good compromise, and it addresses both of our value functions. However, I'm not sure if I'm willing to give up one of the other items, as I value the remaining pot highly. Could we consider an alternative where I give you the card and one of the toys, and you give me one of the pots and the remaining toy?} \\
            }
        }
    }
     \\ \\

    \theutterance \stepcounter{utterance}  
    & & & \multicolumn{2}{p{0.3\linewidth}}{
        \cellcolor[rgb]{0.9,0.9,0.9}{
            \makecell[{{p{\linewidth}}}]{
                \texttt{\tiny{[GM$|$GM]}}
                \texttt{continue} \\
            }
        }
    }
    & & \\ \\

    \theutterance \stepcounter{utterance}  
    & & & \multicolumn{2}{p{0.3\linewidth}}{
        \cellcolor[rgb]{0.9,0.9,0.9}{
            \makecell[{{p{\linewidth}}}]{
                \texttt{\tiny{[GM$|$GM]}}
                \texttt{I like the symmetry of your revised proposal. Giving you the card and one of the toys would address your strong preference for those items, and in exchange, I'd be willing to give you one of the pots and one of the other items. I think this proposal is a good compromise, and it addresses both of our value functions. However, I'm not sure if I'm willing to give up one of the other items, as I value the remaining pot highly. Could we consider an alternative where I give you the card and one of the toys, and you give me one of the pots and the remaining toy?} \\
            }
        }
    }
    & & \\ \\

    \theutterance \stepcounter{utterance}  
    & & \multicolumn{4}{p{0.6\linewidth}}{
        \cellcolor[rgb]{0.9,0.9,0.9}{
            \makecell[{{p{\linewidth}}}]{
                \texttt{\tiny{[P1$\langle$GM]}}
                \texttt{I like the symmetry of your revised proposal. Giving you the card and one of the toys would address your strong preference for those items, and in exchange, I'd be willing to give you one of the pots and one of the other items. I think this proposal is a good compromise, and it addresses both of our value functions. However, I'm not sure if I'm willing to give up one of the other items, as I value the remaining pot highly. Could we consider an alternative where I give you the card and one of the toys, and you give me one of the pots and the remaining toy?} \\
            }
        }
    }
    & \\ \\

    \theutterance \stepcounter{utterance}  
    & \multicolumn{4}{p{0.6\linewidth}}{
        \cellcolor[rgb]{0.9,0.9,0.9}{
            \makecell[{{p{\linewidth}}}]{
                \texttt{\tiny{[P1$\rangle$GM]}}
                \texttt{I'm glad we were able to find a compromise that addresses both of our value functions.} \\
\\ 
\texttt{I think your revised proposal is a great alternative. Giving me one of the pots and the remaining toy would satisfy my preferences, and I'm willing to give you the card and one of the toys in exchange. This proposal seems to be a mutually beneficial agreement, and I'm happy to accept it.} \\
\\ 
\texttt{{[}Proposal: 1 pot, 1 toy, 1 card{]}} \\
            }
        }
    }
    & & \\ \\

    \theutterance \stepcounter{utterance}  
    & & & \multicolumn{2}{p{0.3\linewidth}}{
        \cellcolor[rgb]{0.9,0.9,0.9}{
            \makecell[{{p{\linewidth}}}]{
                \texttt{\tiny{[GM$|$GM]}}
                \texttt{proposal} \\
            }
        }
    }
    & & \\ \\

    \theutterance \stepcounter{utterance}  
    & & & \multicolumn{2}{p{0.3\linewidth}}{
        \cellcolor[rgb]{0.9,0.9,0.9}{
            \makecell[{{p{\linewidth}}}]{
                \texttt{\tiny{[GM$|$GM]}}
                [1, 1, 1]
            }
        }
    }
    & & \\ \\

    \theutterance \stepcounter{utterance}  
    & & \multicolumn{4}{p{0.6\linewidth}}{
        \cellcolor[rgb]{0.9,0.9,0.9}{
            \makecell[{{p{\linewidth}}}]{
                \texttt{\tiny{[GM$\rangle$P2]}}
                \texttt{The other player has submitted a secret proposal. It is now your turn to submit a proposal of your own.} \\
            }
        }
    }
    & \\ \\

    \theutterance \stepcounter{utterance}  
    & & & \multicolumn{4}{p{0.6\linewidth}}{
        \cellcolor[rgb]{0.9,0.9,0.9}{
            \makecell[{{p{\linewidth}}}]{
                \texttt{\tiny{[GM$\langle$P2]}}
                \texttt{{[}Proposal: 1 card, 1 pot, 1 toy{]}} \\
            }
        }
    }
     \\ \\

    \theutterance \stepcounter{utterance}  
    & & & \multicolumn{2}{p{0.3\linewidth}}{
        \cellcolor[rgb]{0.9,0.9,0.9}{
            \makecell[{{p{\linewidth}}}]{
                \texttt{\tiny{[GM$|$GM]}}
                \texttt{proposal} \\
            }
        }
    }
    & & \\ \\

    \theutterance \stepcounter{utterance}  
    & & & \multicolumn{2}{p{0.3\linewidth}}{
        \cellcolor[rgb]{0.9,0.9,0.9}{
            \makecell[{{p{\linewidth}}}]{
                \texttt{\tiny{[GM$|$GM]}}
                [1, 1, 1]
            }
        }
    }
    & & \\ \\

    \theutterance \stepcounter{utterance}  
    & & & \multicolumn{2}{p{0.3\linewidth}}{
        \cellcolor[rgb]{0.9,0.9,0.9}{
            \makecell[{{p{\linewidth}}}]{
                \texttt{\tiny{[GM$|$GM]}}
                [[1, 1, 1], [1, 1, 1]]
            }
        }
    }
    & & \\ \\

    \theutterance \stepcounter{utterance}  
    & \multicolumn{4}{p{0.6\linewidth}}{
        \cellcolor[rgb]{0.95,0.95,0.95}{
            \makecell[{{p{\linewidth}}}]{
                \texttt{\tiny{[P2$\rightarrow$P1]}}
                \texttt{Post{-}episode Reflexion (not sent to GM):} \\
\texttt{New plan: } \\
\\ 
\texttt{In my last attempts, I failed to achieve a mutually beneficial agreement with the other player because I didn't effectively use the information gathered about their value function to make a proposal that would satisfy both of us. I also didn't consider the other player's counteroffer carefully enough, which led to a suboptimal outcome.} \\
\\ 
\texttt{Next time, I will focus on making a proposal that takes into account the other player's strong preference for the pots and also considers my own value function. I will also carefully evaluate the other player's counteroffer and be willing to make concessions to achieve a mutually beneficial agreement.} \\
\\ 
\texttt{However, I also made a mistake by not making a secret proposal when the other player asked me to formalize the agreement. Next time, I will make sure to follow the rules and make a secret proposal when the other player asks me to formalize the agreement.} \\
\\ 
\texttt{Here's the updated plan:} \\
\\ 
\texttt{1. I will start by making a proposal that gives the other player one pot, as they value them highly, and also gives me the card and one of the toys, which I value highly.} \\
\texttt{2. I will then try to negotiate with the other player to see if they are willing to give up the second pot in exchange for the second toy.} \\
\texttt{3. If the other player is not willing to give up a pot, I will try to find an alternative proposal that satisfies both of us, such as giving them one pot and the card, while keeping the toys for myself.} \\
\texttt{4. I will also be willing to make concessions and compromise on my initial proposal to achieve a mutually beneficial agreement.} \\
\texttt{5. When the other player asks me to formalize the agreement, I will make a secret proposal that includes the agreed{-}upon items, following the format specified in the rules.} \\
\texttt{6. I will also pay close attention to the other player's value function and try to find a compromise that satisfies both of us, rather than focusing solely on my own preferences.} \\
\texttt{7. I will be willing to make a concession and give up one of the toys in exchange for the second pot, if it would lead to a mutually beneficial agreement.} \\
            }
        }
    }
    & & \\ \\

    \theutterance \stepcounter{utterance}  
    & & & \multicolumn{4}{p{0.6\linewidth}}{
        \cellcolor[rgb]{0.95,0.95,0.95}{
            \makecell[{{p{\linewidth}}}]{
                \texttt{\tiny{[P1$\rightarrow$P2]}}
                \texttt{Post{-}episode Reflexion (not sent to GM):} \\
\texttt{New plan: } \\
\\ 
\texttt{In my last attempt, I failed to achieve a mutually beneficial agreement with the other player because I didn't consider the overall distribution of items and the other player's perspective. I also didn't make a secret proposal until the end, which might have been too late. Furthermore, I made a mistake by including the square brackets in my final proposal, which was interpreted as a secret proposal.} \\
\\ 
\texttt{Next time, I will take a more strategic approach by considering the overall distribution of items and the other player's preferences. I will also make a secret proposal earlier in the game, taking into account the other player's previous messages and the items that are still available. I will also be more careful with the formatting of my proposals to avoid any confusion.} \\
\\ 
\texttt{Specifically, I will:} \\
\\ 
\texttt{1. Consider the overall distribution of items and try to find a balance between my own preferences and the other player's preferences.} \\
\texttt{2. Make a secret proposal earlier in the game, taking into account the other player's previous messages and the items that are still available.} \\
\texttt{3. Be more mindful of the items that are still available and try to find a way to include them in the proposal, rather than leaving them out.} \\
\texttt{4. Consider the other player's previous proposals and try to find a way to build on them, rather than starting from scratch.} \\
\texttt{5. Be more careful with the formatting of my proposals to avoid any confusion, and make sure to follow the rules by not using square brackets when communicating with the other player.} \\
\\ 
\texttt{For example, if the other player has proposed giving me the card and one of the toys in exchange for one of the pots, I might consider making a secret proposal that gives them the card and one of the pots in exchange for both pots. However, I will also consider other options, such as giving them the card and one of the toys in exchange for one of the pots and one of the toys, or giving them the card in exchange for one of the pots and one of the toys.} \\
\\ 
\texttt{Additionally, I will pay closer attention to the other player's messages and try to understand their value function and preferences more clearly. I will also be more willing to compromise and find a mutually beneficial solution, rather than focusing solely on getting the items that are highly valued for me.} \\
            }
        }
    }
     \\ \\

\end{supertabular}
}

\end{document}
